% Shared preamble for every reading edition in this project.
% A book's main.tex defines \BookLang and \FrankLib (and \BookGloss, if the
% glosses are not in English), then \BookTitle, \BookAuthor, \BookTitleLatin
% and \BookAuthorLatin, then inputs this file.
% Nothing here names a book, and nothing here names a language either: what a
% language decides lives in lib/lang/<code>.tex (docs/languages.md section 6).
% =====================================================================
%  THE ILYA FRANK READING EDITIONS  —  the preamble every book shares
%  A reading edition after the method of Ilya Frank.
%
%  ENGINE:  lualatex main.tex, via ./build.sh   (run twice for the ToC)
%
%  Every passage is given up to five times:
%    1. the whole sentence with its help written in -- the short vowels for
%       Persian and Arabic, the reading over each word (or each chunk) for
%       Japanese and Chinese
%    2. the reading alone, chunk after chunk: kana, or pinyin (only where
%       the chunks carry words: Japanese, Chinese)
%    3. broken into small chunks, each beside a gloss
%    4. the same text bare, the way the language is really printed
%       (only where that differs from 1: Persian, Arabic, Japanese, Chinese)
%    5. the same text once more in the language's other form: nastaliq for
%       Persian, vertical for Japanese and Chinese; nothing for Arabic or Italian
%  A language numbers only the passes it has (the registry's labels): Persian
%  reads 1 the attempt, 2 the chunks, 3 bare, 4 nastaliq.
%
%  The text is written ONCE, in the \ch (\chr, \chw, \chrw) chunks.  Every
%  other pass is assembled from those same tokens, so they can never drift apart.
%
%  WHAT A LANGUAGE FILE DEFINES (lib/lang/<code>.tex), and this file uses:
%    \FrankLangName            babel's name for the language, with its
%                              \babelprovide and \babelfont already issued
%    \FrankRTLtrue/false       which side the chunk sits on; bidi isolates
%    \FrankHasBaretrue/false   whether the bare pass exists
%    \FrankHasAlttrue/false    whether \FrankAltFont exists (a font switch:
%                              the title, the chapter number, and the last
%                              pass unless the language has a vertical pass)
%    \FrankHasVerttrue/false   the last pass is \FrankVertPass (defined by the file)
%    \FrankHasReadingtrue/false  \chr's kana is set (over the chunk, and as
%                              the first line of the gloss); \jruby exists
%    \FrankHasAloudtrue/false  whether the reading pass exists: every
%                              chunk's reading alone, right after pass 1
%                              (a language whose chunks carry words)
%    frank_strip(s)            Lua: the bare form of a chunk (marks dropped)
%    frank_latin(s)            Lua: the label's digits as ASCII, for the ids
%    \FrankHowTo               the "How to read this" text of the front matter
%  and may define \FrankChunkSep (the seam between two chunks in the rebuilt
%  passes; a space unless the file says otherwise), \jruby{text}{kana},
%  \FrankRuby{word}{reading} (the reading over one word of a \chw or \chrw;
%  \jruby unless the file says otherwise), \FrankAloudText{reading} (how a
%  reading is set in the reading pass), and \FrankVbPres / \FrankVbPast (the
%  labels of \vb's two extra forms, each printed only where its form is given).
% =====================================================================
\documentclass[11pt]{book}

% A main.tex written before languages were declared defines neither of these;
% it is a Persian book that finds this file two directories up.
\providecommand{\BookLang}{fa}
\providecommand{\FrankLib}{../../lib}
% The language the glosses are WRITTEN in, which is not always the language
% being taught: an Italian learning English reads English chunks glossed in
% Italian.  Absent it is English, so every book written before the macro
% existed means exactly what it meant (docs/languages.md sections 3 and 11).
\providecommand{\BookGloss}{en}

% ─────────── PAGE + SIZE KNOBS (tune these first) ───────────
\usepackage[paperwidth=105mm,paperheight=140mm,
            top=9mm,bottom=11mm,inner=9mm,outer=9mm]{geometry}
\newcommand{\SzChunk}{\fontsize{13}{25}\selectfont}  % pass 2, the annotated text
\newcommand{\SzGloss}{\fontsize{8.2}{11}\selectfont} % transliteration (and the kana line)
\newcommand{\SzDict} {\fontsize{7.6}{10}\selectfont} % dictionary / vocabulary forms
\newcommand{\SzTrans}{\fontsize{7.2}{9.4}\selectfont}% the meaning (and the ruby kana)
\newcommand{\SzPlain}{\fontsize{15}{27}\selectfont}  % passes 1 and 3, no glosses
\newcommand{\SzNast} {\fontsize{15}{36}\selectfont}  % pass 4 in the alternate face (nastaliq needs air)
\newcommand{\SzToc}  {\fontsize{7.4}{9.4}\selectfont}  % the printed table of contents
% ─────────────────────────────────────────────────────────────

\usepackage{fontspec}
\usepackage[bidi=basic,english,provide+=*]{babel}
% The document's own language -- what babel was just loaded with, and what the
% roman below is set for.  It is not the gloss language (\FrankGlossName, and
% a book may be glossed in Italian) and it is not the book's: it is the name
% of the face and the script the series is set in, which is what a
% transliteration goes back into wherever the glosses around it are set in
% another (\FrankRoman below).
\newcommand{\FrankRomanLang}{english}
% Glyphs the roman lacks come from a fallback face, glyph by glyph: the
% IPA-lite marks an Italian pronunciation line uses (ʎ ɲ ʃ ṡ), the dotted
% letters of a scholarly transliteration.  luaotfload does the substitution
% only for characters TeX Gyre Pagella has no glyph for, so a book whose
% glosses stay inside Pagella's repertoire -- every Persian one -- is set
% exactly as before.  Without DejaVu Serif on the machine the fallback is
% simply not declared and a missing glyph is reported in the log as always.
\IfFontExistsTF{DejaVu Serif}{%
  \directlua{luaotfload.add_fallback("frankfallback", {"DejaVu Serif:mode=node;script=latn;"})}%
  \babelfont{rm}[RawFeature={fallback=frankfallback}]{TeX Gyre Pagella}%
}{%
  \babelfont{rm}{TeX Gyre Pagella}%
}
\babelfont{sf}{TeX Gyre Heros}
\newfontfamily\symbolfont{DejaVu Sans}   % carries the sun the 1951 edition uses
\usepackage{etoolbox}

% ---------------------------------------------------------------------
%  THE LANGUAGE
% ---------------------------------------------------------------------
% The switches the language file sets.  All false until it speaks.
\newif\ifFrankRTL
\newif\ifFrankHasBare
\newif\ifFrankHasAlt
\newif\ifFrankHasVert
\newif\ifFrankHasReading
\newif\ifFrankHasAloud

% The fonts a language is set in are the registry's, lib/languages.json --
% no .tex file names a font (docs/languages.md section 1).  LuaTeX reads the
% registry itself: lualibs, which luaotfload loads, brings utilities.json.
% The file is looked for beside this one, through \FrankLib, from the book
% directory lualatex runs in.  A registry that cannot be read is a real error
% (tex.error gives a "!" line, which build.sh stops on), not a quiet fallback:
% the book would otherwise come out in whatever face fontspec picked.
% No #, %, ~ or literal backslash in a \directlua body: TeX reads it first,
% and ~ is an active character that expands to \nobreakspace -- so Lua's
% "not equal" is written not (a == b) here (see frank_breaks below too).
\directlua{
  frank_registry = nil
  function frank_reg_load()
    if frank_registry then return frank_registry end
    local path = "\FrankLib/languages.json"
    local f = io.open(path, "r")
    if not f then
      tex.error("frank: cannot read the language registry at " .. path)
      frank_registry = {} ; return frank_registry
    end
    local s = f:read("*a") ; f:close()
    if not (utilities and utilities.json and utilities.json.tolua) then
      tex.error("frank: this LuaTeX has no JSON reader (lualibs); cannot read " .. path)
      frank_registry = {} ; return frank_registry
    end
    frank_registry = utilities.json.tolua(s) or {}
    return frank_registry
  end
  function frank_reg(lang, ...)
    local t = frank_reg_load()[lang]
    for _, k in ipairs({...}) do
      if not (type(t) == "table") then return "" end
      t = t[k]
    end
    if type(t) == "table" then return table.concat(t, ",") end
    if t == nil or t == false then return "" end
    return tostring(t)
  end
  function frank_reg_fonts(lang, key, fbkey)
    local out, n = {}, 0
    local function add(v)
      if type(v) == "table" then
        for _, x in ipairs(v) do add(x) end
      elseif type(v) == "string" and string.len(v) > 0 then
        n = n + 1 ; out[n] = v
      end
    end
    local tex_rec = frank_reg_load()[lang]
    tex_rec = type(tex_rec) == "table" and tex_rec.tex or nil
    if type(tex_rec) == "table" then
      add(tex_rec[key])
      if fbkey and string.len(fbkey) > 0 then add(tex_rec[fbkey]) end
    end
    return table.concat(out, ",")
  end
}
% \FrankPickFont{\Name}{<code>}{<key>}{<fallbacks key>}
%   \Name becomes the first face, of the registry's tex.<key> followed by
%   tex.<fallbacks key>, that is installed (fontspec's \IfFontExistsTF, which
%   sees OSFONTDIR and so the faces in lib/fonts).  When none is, \Name is the
%   first name anyway and a warning says so: fontspec then fails loudly on it,
%   which is right -- a book without a font for its language must not build.
\makeatletter
\newcommand{\FrankPickFont}[4]{%
  \edef\Frank@cands{\directlua{tex.sprint(-2, frank_reg_fonts("#2","#3","#4"))}}%
  \def#1{}%
  \def\Frank@first{}%
  \renewcommand*{\do}[1]{%
    \ifdefempty{\Frank@first}{\def\Frank@first{##1}}{}%
    \ifdefempty{#1}{\IfFontExistsTF{##1}{\def#1{##1}}{}}{}}%
  \expandafter\docsvlist\expandafter{\Frank@cands}%
  \ifdefempty{#1}{%
    \PackageWarningNoLine{frank-preamble}{none of the registry's fonts for
      #2/#3 is installed (\Frank@cands); trying \Frank@first}%
    \let#1\Frank@first}{}}
\makeatother

\input{\FrankLib/lang/\BookLang.tex}

% What a language file may leave to the defaults:
%   \FrankChunkSep -- the seam between two chunks when a pass is rebuilt from
%   the buffers.  A space, for every language written with spaces.
%   \jruby{text}{kana} -- a language without a reading prints the text alone,
%   so \chr can be written for every language and set the same way.
\providecommand{\FrankChunkSep}{\space}
\providecommand{\jruby}[2]{#1}
%   \FrankRuby{word}{reading} -- the reading over ONE word of a \chw or
%   \chrw.  Japanese's is its \jruby, measured per word as it is per chunk;
%   a language that gives neither prints the word alone.
\providecommand{\FrankRuby}[2]{\jruby{#1}{#2}}
%   \FrankAloudText{reading} -- a chunk's reading in the reading pass, as it
%   stands in the language's face (Chinese puts its pinyin in the roman).
\providecommand{\FrankAloudText}[1]{#1}
%   \FrankVbPres, \FrankVbPast -- the two labels \vb prints before its third
%   and fifth arguments, each only when that argument is not blank (see \vb
%   below).  "pres." and "past" are the Persian stems; a language whose verb
%   is given by other forms renames them (ja: the -masu stem and the -te
%   form; tr: the -iyor present and the aorist; zh: the split and the "can't"
%   form).  Every language file sets them, and the registry's "vb_labels"
%   carries the same pair for the reader, which cannot read this file;
%   lib/newlang.py --check refuses to let the two spellings drift apart.
%   These defaults are what an unmigrated language file would fall back to.
\providecommand{\FrankVbPres}{pres.}
\providecommand{\FrankVbPast}{past}

% ---------------------------------------------------------------------
%  THE GLOSS LANGUAGE
% ---------------------------------------------------------------------
% The book TEACHES \BookLang and its glosses are WRITTEN in \BookGloss.  Those
% are two questions, and until a book could answer the second every gloss,
% vocabulary line and contents row below said `english' outright.
% \FrankGlossName is what they say now.
%
% A gloss language is PROSE and need not be one of the eight this toolbox
% teaches -- somebody may gloss Persian in Spanish, and nobody studies Spanish
% here.  So the name is asked of the registry first, which for a language the
% toolbox does teach also gives the direction it runs in and the face it is
% set in; and of babel's own locale file otherwise, under a name of ours,
% because babel takes any name and reads the locale from `import'.  That is
% why no list of the languages one may write in is kept here either
% (docs/languages.md section 1).
%
% `en' is the one code nothing is provided for: babel was loaded with
% `english' as the document's language at the top of this file, so
% \FrankGlossName is already a language babel has.  A book that says nothing
% -- every book written before \BookGloss existed -- therefore reaches
% \begin{otherlanguage}{english} through the very tokens it always did, which
% is how the Persian edition rebuilds to the same PDF (section 11).
%
% frank_gloss_opts(code) builds the option list of the gloss's \babelprovide.
% `import' is a BCP-47 tag, whose region subtag is written in capitals, and
% babel names its locale files after the tag: en-gb finds babel-en-gb.ini on a
% case-sensitive disk, which is not there, and says nothing about it.
% Then every option the registry gives that language EXCEPT onchar, which
% hands a whole script to ONE babel language: the book's own must keep it, or
% a Persian gloss would take the Arabic of the text it is glossing away from
% the language teaching it.  Nothing is lost by dropping it -- every gloss
% below names its language outright, so none of them needs switching on sight.
% No #, %, ~ or literal backslash in a \directlua body, and no Lua comment
% either: TeX reads the body first and turns its line ends into spaces, so a
% -- comment swallows the rest of the function (that is why every block here
% is commented from the outside).  No Lua character class, for the % in it.
\directlua{
  function frank_gloss_opts(code)
    local tag = frank_reg(code, "tex", "import")
    if string.len(tag) == 0 then tag = code end
    local i = string.find(tag, "-", 1, true)
    if i then tag = string.sub(tag, 1, i) .. string.upper(string.sub(tag, i + 1)) end
    local out = "import=" .. tag
    for piece in string.gmatch(frank_reg(code, "tex", "options") .. ",", "([^,]*),") do
      if string.len(piece) > 0 and not (string.sub(piece, 1, 7) == "onchar=") then
        out = out .. "," .. piece
      end
    end
    return out
  end
}
\makeatletter
\edef\FrankGlossName{\directlua{tex.sprint(-2, frank_reg("\BookGloss", "tex", "babel"))}}
\ifdefempty{\FrankGlossName}{\def\FrankGlossName{frankgloss}}{}
% Which way the glosses run.  The registry answers for the eight; a code it
% does not carry is one of the prose languages the studio writes notes in
% (docs/languages.md section 5), every one of them Latin-script and left to
% right, and guessing otherwise would set a page backwards without a word.
\newif\ifFrankGlossRTL
\edef\Frank@gdir{\directlua{tex.sprint(-2, frank_reg("\BookGloss", "dir"))}}
\ifdefstring{\Frank@gdir}{rtl}{\FrankGlossRTLtrue}{}
% \FrankRoman{...} is what in a gloss is Latin script BY NATURE and not by
% choice: the transliteration line, the sound beside each word of a vocabulary
% entry, and the two labels a \vb prints (every language's pair is a Latin
% abbreviation -- pres., masdar, -te).  Where the gloss language brings a face
% of its own -- Persian, Arabic, Japanese -- that face would set them, and it
% has no italic (babel substitutes the upright and says so once in the log)
% and shapes their diacritics with the wrong script: risālatan ṭawīlatan came
% out of Vazirmatn with the dot of the ṭ on the letter before it.  So they go
% back into the language the roman is set for, which is also what makes each
% of them a left-to-right run of its own.  For a gloss the roman already sets
% -- every Latin-script one, English included -- this is nothing at all, and
% the Persian edition's vocabulary lines keep the tokens they had (§11).
\newcommand{\FrankRoman}[1]{#1}
% Provide it only where nobody has: `english' is the document's own language,
% and lib/lang/<code>.tex has already provided the gloss's when a book is
% glossed in the language it teaches (an English edition, a monolingual one).
\newif\ifFrank@needgloss  \Frank@needglosstrue
\ifdefstring{\FrankGlossName}{english}{\Frank@needglossfalse}{}
\ifdefstrequal{\FrankGlossName}{\FrankLangName}{\Frank@needglossfalse}{}
\ifFrank@needgloss
  \edef\Frank@gloss@provide{\noexpand\babelprovide[%
    \directlua{tex.sprint(-2, frank_gloss_opts("\BookGloss"))}]{\FrankGlossName}}%
  \Frank@gloss@provide
  % The face, where the registry names one: Persian glosses want Vazirmatn,
  % Italian ones the roman the core already loaded (tex.main is null for every
  % Latin-script language, and a second face would set one alphabet two ways).
  \edef\Frank@gloss@main{\directlua{tex.sprint(-2, frank_reg("\BookGloss", "tex", "main"))}}%
  \ifdefempty{\Frank@gloss@main}{}{%
    \FrankPickFont{\FrankGlossFont}{\BookGloss}{main}{main_fallbacks}%
    \edef\Frank@gloss@font{\noexpand\babelfont[\FrankGlossName]{rm}[%
      Script=\directlua{tex.sprint(-2, frank_reg("\BookGloss", "tex", "script"))},%
      Language=\directlua{tex.sprint(-2, frank_reg("\BookGloss", "tex", "language"))},%
      Renderer=HarfBuzz]{\FrankGlossFont}}%
    \Frank@gloss@font
    \renewcommand{\FrankRoman}[1]{\foreignlanguage{\FrankRomanLang}{#1}}}%
\fi
\makeatother
% Whether a word of the book's language, set inside a gloss, has to be
% isolated from what surrounds it: it does exactly when the two run opposite
% ways (\pw below).  Every book glossed in English asks this of \ifFrankRTL,
% which is what this was until the gloss could run the other way.
\newif\ifFrankMixedDir
\ifFrankRTL
  \ifFrankGlossRTL\else\FrankMixedDirtrue\fi
\else
  \ifFrankGlossRTL\FrankMixedDirtrue\fi
\fi

% One pairing is refused rather than set backwards.  A gloss block is ONE
% paragraph and runs the gloss language's way; every island of the other
% direction in it has to be isolated, and \pw isolates the ones this file
% knows about -- the words of the book's own language.  When the gloss runs
% left to right those are the only islands there are, because a
% transliteration and a label are Latin and Latin is then the paragraph's own
% direction: that is every book in the toolbox today, the Persian one
% included.  When the gloss runs right to left and the book does too, the
% islands are the romanisations, one to an entry, with the language's own
% words on both sides of them, and the bidi algorithm places them correctly
% (an Arabic text glossed in Persian sets exactly as it should).  But when the
% gloss runs right to left and the BOOK runs left to right, a vocabulary line
% is a chain of Latin runs joined by punctuation -- the period of `pres.', the
% semicolons, the middle dots -- and every one of those neutral characters
% takes the paragraph's direction and throws the entries it separates out of
% order.  Isolating each entry instead would put it in an \mbox, which cannot
% break across a line, and a \vb is longer than the gloss column.
\ifFrankGlossRTL\ifFrankRTL\else
  \PackageError{frank-preamble}{%
    \BookGloss\space glosses cannot be written for a \BookLang\space book}{%
    The glosses run right to left and the text they gloss runs left to
    right.\MessageBreak
    A vocabulary line would then come out with its entries in the wrong
    order: the punctuation between them takes the direction of the
    paragraph, not of the words on either side.\MessageBreak
    Gloss this book in a left-to-right language, or give a right-to-left
    gloss to a right-to-left text -- Persian glosses of an Arabic edition
    are set correctly.}%
\fi\fi

\usepackage{xcolor}
\usepackage{environ}   % \BODY: lets a sentence be collected before it is set

% ---------------------------------------------------------------------
%  NAVIGATION -- PDF outline, named destinations, printed contents
% ---------------------------------------------------------------------
% The unit the reader navigates by is the PARAGRAPH, never the sentence:
% one \parend or \chapend per paragraph, so each already opens a page.
%
%   identity      chapter + paragraph, both in LATIN digits
%   destination   par.3.12   (chapter 3, paragraph 12)
%   human label   ۳.۱۲       -- the book's own digits (the language's)
%
% hyperref must come after babel, and bookmark after hyperref.
\usepackage{hyperref}
\hypersetup{unicode,bookmarks,bookmarksnumbered=false,hidelinks,
            bookmarksopen,bookmarksdepth=2,
            pdftitle={\BookTitleLatin},pdfauthor={\BookAuthorLatin}}
\usepackage{bookmark}

% The language's digits are the label, Latin digits are the id, so a
% destination name has to be transliterated: frank_latin, from the language
% file, does it.  \directlua is expandable, which is what lets this survive an
% \edef and be built into a name argument.
\newcommand{\FrankLatin}[1]{%
  \directlua{tex.sprint(-2, frank_latin("\luaescapestring{#1}"))}}

% \parstart is handed chapter.paragraph, because that is what makes a unique
% destination.  But the page itself prints paragraph.subparagraph (NOTES §3),
% so showing "۱.۲" in the contents beside a page headed "۲.۱" reads as the same
% number reversed.  The chapter is already the level above, in the outline and
% under its own heading in the printed list, so the entry shows the paragraph
% alone and the reader sees ۲ against a page of ۲.۱, ۲.۲, ...
% No Lua patterns here: TeX reads the argument of \directlua first, and % is
% still its comment character, so "%.(.*)$" would silently comment out the rest
% of the line and the function would return nothing.  A plain find avoids it.
\directlua{
  function frank_parno(s)
    local i = s:find(".", 1, true)
    if i then return s:sub(i + 1) end
    return s
  end
}
\newcommand{\FrankParNo}[1]{%
  \directlua{tex.sprint(-2, frank_parno("\luaescapestring{#1}"))}}

\makeatletter
% A named destination that is legal in VERTICAL mode.  \Hy@raisedlink drops
% the whatsit straight onto the vertical list when there is no line to hang
% it from, so \parstart contributes no box, no glue and no paragraph --
% \hypertarget would have had to start one.
\newcommand{\FrankAnchor}[1]{%
  \Hy@raisedlink{\hyper@anchorstart{#1}\hyper@anchorend}}

% ---- the two kinds of contents line ----
% \contentsline{frankpar}{\frankparentry{label}{incipit}}{page}{} is what
% \parstart writes; \l@frankpar sets it.  Both are paragraphs of the GLOSS
% language and run its way, so that the label comes before the incipit and
% the page number lands at the end of the line for whoever is reading the
% apparatus -- left to right for an English or Italian gloss, the other way
% for a Persian one.  The label is a number in an RTL run and needs
% the \mbox isolate or ۳.۱۲ comes out ۱۲.۳ (NOTES §9); the incipit is one
% continuous phrase of the language, so it is left a plain run and is free to
% break over two lines.  \hfil between two \parbox es is a legal breakpoint
% under \raggedright, hence the enclosing \hbox to\linewidth (NOTES §9).
% \parbox[t] takes its baseline from the FIRST BOX on its vertical list, so
% both columns must be in horizontal mode before anything else -- \color
% alone leaves a whatsit there, the box height comes out 0, and the page
% number drops half a line below the line it belongs to.
\newcommand{\FrankTocRow}[2]{%
  \par\noindent
  \hbox to\linewidth{%
    \parbox[t]{0.84\linewidth}{\raggedright
      \hyphenpenalty=10000\exhyphenpenalty=10000 \leavevmode#1}%
    \hfil
    \parbox[t]{0.12\linewidth}{\raggedleft\leavevmode\color{faint}#2}}%
  \par}

\newcommand{\frankparentry}[2]{% #1 = ۳.۱۲ (chapter.paragraph), #2 = the incipit
  \edef\Frank@dest{par.\FrankLatin{#1}}%
  \expandafter\hyperlink\expandafter{\Frank@dest}{%
    \mbox{\foreignlanguage{\FrankLangName}{\FrankParNo{#1}}}\hskip0.9em\relax
    \foreignlanguage{\FrankLangName}{#2}}}

\newcommand{\l@frankpar}[2]{%
  \begin{otherlanguage}{\FrankGlossName}%
    {\SzToc\color{ink}\FrankTocRow{#1}{#2}}%
  \end{otherlanguage}%
  \addvspace{0.15ex}}

\newcommand{\frankchapentry}[1]{% #1 = ۳
  \edef\Frank@dest{ch.\FrankLatin{#1}}%
  \expandafter\hyperlink\expandafter{\Frank@dest}{%
    \mbox{\foreignlanguage{\FrankLangName}{#1}}}}

\newcommand{\l@frankchap}[2]{%
  \addvspace{1.8ex}%
  \begin{otherlanguage}{\FrankGlossName}%
    {\normalfont\footnotesize\color{ink}\FrankTocRow{#1}{#2}}%
  \end{otherlanguage}%
  \addvspace{0.9ex}}

% \parstart{۳.۱۲}{first six words of pass 1}
%   The destination, the outline bookmark and the contents line for one
%   paragraph.  It must produce NO visible output and NO vertical space:
%   verify_book.py proves the body still reproduces the source word for
%   word, and the four passes must stay exactly where they were.
%   Every name is expanded HERE -- \bookmark and \addtocontents both write
%   at shipout, by which time \Frank@dest has moved on to another paragraph.
\newcommand{\parstart}[2]{%
  \edef\Frank@dest{par.\FrankLatin{#1}}%
  \FrankAnchor{\Frank@dest}%
  \edef\Frank@bkm{\noexpand\bookmark[dest=\Frank@dest,level=1]}%
  % \bookmark writes its text out to the .out file, and a \directlua sitting in
  % that text is gone by the time it is written -- the number simply vanished
  % and the title began with the dash.  Expand it here, like the destination.
  \edef\Frank@no{\FrankParNo{#1}}%
  \expandafter\Frank@bkm\expandafter{\Frank@no\space— #2}%
  \addtocontents{toc}{%
    \protect\contentsline{frankpar}{\protect\frankparentry{#1}{#2}}{\thepage}{}}}

% the same, for a chapter: level 0, and its own contents line
\newcommand{\chapstart}[1]{%
  \edef\Frank@dest{ch.\FrankLatin{#1}}%
  \FrankAnchor{\Frank@dest}%
  \edef\Frank@bkm{\noexpand\bookmark[dest=\Frank@dest,level=0]}%
  \Frank@bkm{#1}%
  \addtocontents{toc}{%
    \protect\contentsline{frankchap}{\protect\frankchapentry{#1}}{\thepage}{}}}

% The printed contents, called once by the front matter.  The front matter
% is numbered in roman and the body restarts at 1, so however many pages
% this list runs to it cannot move the numbers printed inside it -- which
% is what lets build.sh's two lualatex passes be enough.
\newcommand{\frankcontents}{%
  \thispagestyle{fancy}%
  \FrankAnchor{frank.toc}%
  \bookmark[dest=frank.toc,level=0]{Contents}%
  {\centering\normalfont\small\textbf{Contents}\par}%
  \vspace{2.5mm}%
  \@starttoc{toc}%
  \par}
\makeatother

% Pass 1 carries the help -- the short vowels, the way a children's edition
% does.  Passes 3 and 4 must not, but the text is only written once -- so the
% marks are filtered out of the copy that feeds them, by the language file's
% frank_strip (the harakat for Persian and Arabic; nothing for the others).
\newcommand{\Strip}[1]{\directlua{tex.sprint(frank_strip("\luaescapestring{#1}"))}}

% The 1951 edition writes a colon with no space after it -- بکند:بوی,
% فاسقهایی:سیرابی‌فروش.  Arabic script offers no break inside a joined run, so
% such a pair is one unbreakable token; at 13pt in the 0.34\linewidth chunk
% column فاسِقهایی:سیرابی‌فُروش، came out 41pt too wide and ran off the page.
% Re-chunking cannot help -- the token itself is too wide -- and splitting it
% would put a space where the source has none and fail the fidelity test.  So
% the colon is given a break opportunity, in the chunk column only: passes 1, 3
% and 4 set the sentence across the full measure, where it already fits.
%
% The backslash is built from its char code.  Writing it literally does not
% work: TeX reads the body of \directlua before Lua ever sees it, so a \\ in
% there is an undefined control sequence, not an escaped backslash.
%
% A pre_linebreak_filter callback was tried first and never fired on this text:
% babel's bidi=basic boxes the run, so the glyphs are not on the list the
% callback is handed.
% The same problem arrives twice more.  The edition marks a pause with a row of
% 58 dots (chapter 3 paragraphs 73 and 90), which is one unbreakable token 204pt
% wider than the chunk column, and it writes «مرگ... مرگ...» with ellipses inside
% a chunk.  A period gets the same treatment as a colon: elsewhere a period is
% followed by a space, so nothing else moves.
% No Lua patterns: TeX reads the body of \directlua first, so a %. in there is
% a comment, not an escaped dot.  Walking the bytes avoids the question, and is
% safe over UTF-8 because a break is only ever appended after an ASCII byte.
% And no # either -- string.len(s), never #s.  TeX DOUBLES a macro parameter
% character on its way into \directlua, so #s reaches Lua as ##s, which is the
% length of a length: "attempt to get length of a number value".  \directlua
% then abandons the chunk, tex.sprint never runs, and the Persian column comes
% out EMPTY -- silently, because that error does not start with a ! and so read
% as a clean build.  It cost the book every chunk in passes 2 and 3 once.
\directlua{
  function frank_breaks(s)
    local br = string.char(92) .. "allowbreak "
    local out, n = {}, 0
    for i = 1, string.len(s) do
      local c = s:sub(i, i)
      n = n + 1 ; out[n] = c
      if c == ":" or c == "." then n = n + 1 ; out[n] = br end
    end
    return table.concat(out)
  end
}
\newcommand{\FrankBreaks}[1]{%
  \directlua{tex.sprint(frank_breaks("\luaescapestring{#1}"))}}
\usepackage{fancyhdr}

% Kindle is greyscale: everything below is a grey, not a colour.
\definecolor{ink}  {gray}{0.10}
\definecolor{gloss}{gray}{0.42}
\definecolor{tlit} {gray}{0.22}
\definecolor{dict} {gray}{0.30}
\definecolor{faint}{gray}{0.70}

\pagestyle{fancy}\fancyhf{}
\renewcommand{\headrulewidth}{0pt}
\fancyfoot[C]{\color{faint}\normalfont\footnotesize\thepage}

% ---------------------------------------------------------------------
%  THE PASSES
% ---------------------------------------------------------------------
% \ch{colour}{text}{transliteration}{vocabulary forms}{meaning}
%   The text is written ONCE, with its help (the vowels).  \ch typesets the
%   chunk with its gloss and appends a bare copy to \FrankBuf, out of which
%   passes 3 and 4 are rebuilt, so the passes cannot drift apart.
% \chr{colour}{text}{kana}{transliteration}{vocabulary forms}{meaning}
%   The same for a language with a reading (docs/languages.md section 3): the
%   kana is set over the chunk in passes 1 and 2 and heads the gloss.  A
%   distinct macro, not a sixth argument, so every parser knows the arity
%   from the name.  In a language without a reading it behaves as \ch.
% \chw{colour}{text}{transliteration}{vocabulary forms}{meaning}{words}
% \chrw{colour}{text}{kana}{transliteration}{vocabulary forms}{meaning}{words}
%   \ch and \chr carrying the chunk's WORD LINE, `山(やま) へ 柴刈り(しばかり) に 、'
%   (lib/wordline.py): pass 1 and the chunk column set the text a word at a
%   time, each word's reading over it.  The words come LAST, so every reader
%   of the older two keeps its argument numbers.
%
% LAYOUT: \FrankFlowtrue puts the glosses inline, the way Frank's printed
%   editions do.  It needs a wide measure -- on a page this narrow a long
%   chunk and its gloss will not fit on one line, and an inline gloss that
%   breaks gets its halves thrown to opposite corners by the bidi algorithm.
%   \FrankFlowfalse gives each chunk its own line, text on the language's
%   side, gloss on the other.
\newif\ifFrankFlow
\FrankFlowfalse

% ------------------- the highlight palette ---------------------------
% Four colours to mark a chunk with, put in \ch's first argument.  They reach
% pass 1 only -- the whole sentence the reader attempts before any help -- so
% a mark can point at a word without giving the game away in the chunks, the
% bare text or the fourth pass, which stay in \c@ink like everything else.
%
% The rest of this book is deliberately greyscale, because it is read on
% e-ink (see NOTES section 3).  On such a screen these will come out as four
% greys, and dark red against green may not separate; on paper or a colour
% screen they will.
\definecolor{frankred}   {HTML}{8E1B1B}
\definecolor{frankblue}  {HTML}{1F4E79}
\definecolor{frankorange}{HTML}{BF5B04}
\definecolor{frankgreen} {HTML}{22643A}
\newcommand{\Cred}   {\color{frankred}}
\newcommand{\Cblue}  {\color{frankblue}}
\newcommand{\Corange}{\color{frankorange}}
\newcommand{\Cgreen} {\color{frankgreen}}

\newcommand{\ChunkCol}{0.34\linewidth}
\newcommand{\GlossCol}{0.62\linewidth}

% \FrankVocBuf keeps the help (pass 1), \FrankBuf drops it (the bare pass and
% the alternate), \FrankAloudBuf holds each chunk's reading alone (the reading
% pass, where the language has one)
\newif\ifFrankCollect
\newcommand{\FrankBuf}{}
\newcommand{\FrankVocBuf}{}
\newcommand{\FrankAloudBuf}{}
\newcommand{\FrankBufReset}{\gdef\FrankBuf{}\gdef\FrankVocBuf{}\gdef\FrankAloudBuf{}}

% ------------------ pieces of the vocabulary line --------------------
% A word of the book's language sitting inside a run of the gloss's.  When the
% two run opposite ways it has to be isolated: two such words separated only
% by a space merge into one run of their own direction and come out in reverse
% order, and \babelsublr+\mbox keeps each one where it was typed.  When they
% run the same way -- an Italian book glossed in Italian, an Arabic one
% glossed in Persian -- only the \mbox is wanted, and \babelsublr would be
% wrong: it opens a run of the direction that is already current, which babel
% does not expect.  For every book glossed in English this is \ifFrankRTL,
% which is the test that stood here while English was the only gloss there was.
\newcommand{\pw}[1]{%
  \ifFrankMixedDir
    \babelsublr{\mbox{\foreignlanguage{\FrankLangName}{\PwScale #1}}}%
  \else
    \mbox{\foreignlanguage{\FrankLangName}{\PwScale #1}}%
  \fi}
\newcommand{\PwScale}{\fontsize{8.6}{10}\selectfont}

% \FrankSound{sound}: the space before a word's sound, and the sound -- or,
% when the sound is blank, NOTHING, not even the space.  A sound slot is left
% blank wherever the conventions give none: in every \vb docs/lang/tr.md,
% es.md and de.md model and in most of it.md's, in the \dw of a Spanish or
% German noun (de.md's own model is \dw{der Tisch, -e}{} table), in a German
% \bw.  Written out as `\ \FrankRoman{\textit{#2}}', the space stayed when
% the sound was not there and met the space the next piece brings, so each
% blank sound set two: in the Turkish fixture every dot of "yağmak · pres.
% yağıyor · aor. yağar · to fall" had 3.79pt before it where a \vb with its
% sounds has 1.89pt, and \dw{vez}{} time and \bw{Werk}{}{work} had the same
% 3.79pt before the next word.  The reader never printed the second space
% (texparse drops an empty sound run), so the PDF now sets the line as the
% reader does.  A filled sound is the very same tokens it always was: nine of
% the eleven fixture editions come out of pdftotext byte for byte as before
% (mini-de among them, whose \bw{Werk}{}{work} closed up without moving a
% line break), and in the other two, mini-tr and mini-es, only line breaks
% move -- the same words on every page of mini-tr, and in mini-es a line of
% pass 1 that now fits at the foot of page 8.
\newcommand{\FrankSound}[1]{\ifblank{#1}{}{\ \FrankRoman{\textit{#1}}}}
\newcommand{\dw}[2]{\pw{#1}\FrankSound{#2}}     % headword + sound
\newcommand{\bw}[3]{\ \textperiodcentered\ \pw{#1}\FrankSound{#2} `#3'} % base of a compound verb
% \vb{form1}{sound1}{form2}{sound2}{form3}{sound3}{meaning}: the same seven
% arguments in every language, the three forms being the language's own
% (lib/lang/<code>.tex names them, the registry's "vb_forms" says them in
% words, \FrankVbPres / \FrankVbPast label the second and the third).
%
% A PAIR WHOSE FORM IS BLANK IS NOT PRINTED -- neither the form nor its label
% nor the dot before it -- and the same test is made on the meaning, which
% always had it.  Two languages need a pair left out, and before the test
% they could only print a label standing over nothing.  Chinese gives a
% separable verb its split and a verb with a complement its "can't" form,
% never both: \vb{睡觉}{shuìjiào}{睡了觉}{shuìle jiào}{}{}{to sleep} came out
% as "... · split 睡了觉 shuìle jiào · can't · to sleep".  And an Arabic verb
% whose masdar is not in use would have printed "· masdar" and nothing.
% Only the FORM decides whether a pair is printed: a sound slot is empty in
% every \vb docs/lang/tr.md, es.md and de.md model and in most of it.md's,
% and a pair with a form and no sound is a pair.  The sound decides only
% whether there is a space for it: each of the three goes through \FrankSound
% above, so \vb{yağmak}{}{yağıyor}{}{yağar}{}{to fall} sets one space before
% each dot and not two.
% \ifblank (etoolbox, loaded above) takes spaces for blank and expands
% nothing in the test, so \vb{a}{}{ }{}{c}{}{} leaves the second pair out too.
%
% Why not an optional argument, or a macro per language, for the pieces a
% language adds (an auxiliary, a verb class, a governed case)?  Both were
% tried: every tool that reads a \vb -- texparse, check_batch and texwrite
% through it, normalize_batch, chunkdiv, assemble -- would have had to learn
% the new shape on the same day, and one that did not failed loudly or, worse,
% quietly: texparse.parse_voc died on \vb[aux. sein]{...} with "expected a
% group at 3", which takes the reader's build down with it, while check_batch
% counted its arguments as 0 and a per-language \vbde was refused as a macro
% not allowed.  So the extras stay where docs/lang/<code>.md already put them,
% in ONE parenthesis after the meaning, items parted by "; ":
% \vb{fahren}{}{fuhr}{}{gefahren}{}{to drive (er fährt; aux. sein)}.
%
% Every pair in every book written before the test is filled, so they print
% exactly as they did: the pdftotext of six fixture editions (mini-fa with
% its 136 \vb, mini-ar, -de, -ja, -it, -en) is byte for byte what it was.
% lib/texparse.py's _voc_runs makes the same test for the reader, and must
% keep making it: a pair the PDF leaves out and the reader prints is one book
% teaching two things.
\newcommand{\vb}[7]{%
  \pw{#1}\FrankSound{#2}%
  \ifblank{#3}{}{\ \textperiodcentered\ \FrankRoman{\FrankVbPres}~\pw{#3}\FrankSound{#4}}%
  \ifblank{#5}{}{\ \textperiodcentered\ \FrankRoman{\FrankVbPast}~\pw{#5}\FrankSound{#6}}%
  \ifblank{#7}{}{\ \textperiodcentered\ #7}}

% ------------------------ the gloss block ----------------------------
% The block is set in the GLOSS language, whatever the book teaches: which
% way it runs, which face sets it, and its locale's own punctuation -- a
% French gloss gets the space before ; : ! ? and the right apostrophe from
% babel rather than from whoever typed it.  Not its break points: the line
% below forbids hyphenation in this column altogether, and always has.
% \glbodyr{kana}{tr}{voc}{en}: the kana line comes first, upright, in the
% language's own font, and only for a language with a reading -- for the
% others the first argument is ignored, so \glbody is \glbodyr with it empty.
% Every line but the meaning is guarded: the transliteration is optional
% wherever the registry's require_tr is false (Turkish gives one only for the
% three cases its conventions name), and an unguarded \textit{} would print a
% blank line the gloss has to be read across.
\newcommand{\glbodyr}[4]{%
  \begin{otherlanguage}{\FrankGlossName}%
    \raggedright\hyphenpenalty=10000\exhyphenpenalty=10000
    \emergencystretch=1em
    \ifFrankHasReading
      \ifblank{#1}{}{{\SzGloss\color{tlit}\foreignlanguage{\FrankLangName}{#1}\par}}%
    \fi
    \ifblank{#2}{}{{\SzGloss\color{tlit}\FrankRoman{\textit{#2}}\par}}%
    \ifblank{#3}{}{{\SzDict\color{dict}#3\par}}%
    {\SzTrans\color{gloss}#4\par}%
  \end{otherlanguage}}
\newcommand{\glbody}[3]{\glbodyr{}{#1}{#2}{#3}}

% The chunk and its gloss side by side.  The two \parbox es are written chunk
% first, and the paragraph's direction puts them on the right side: inside an
% RTL language an \hbox lays its contents out right to left, so the chunk is
% on the RIGHT and the gloss on the left; inside an LTR language the same
% \hbox puts the chunk on the LEFT and the gloss on the right.  No swap is
% needed, and none must be written -- swapping for LTR would put the gloss
% first.  \hfil between the boxes is a legal breakpoint under \raggedright,
% hence the enclosing \hbox to\linewidth (NOTES §9).
\newcommand{\FrankChunkRow}[2]{%
  \par\noindent
  \hbox to\linewidth{%
    \parbox[t]{\ChunkCol}{\raggedright #1}%
    \hfil
    \parbox[t]{\GlossCol}{#2}}%
  \par\addvspace{1.1ex}}

% \ch{colour}{fa}{tr}{voc}{en}
% The first argument is normally empty.  Give it \Cred, \Cblue, \Corange or
% \Cgreen and that chunk is coloured IN PASS 1 ONLY; every other pass takes
% the text from a separate buffer that never sees it.
\newcommand{\FrankVocAppend}[2]{%
  \ifblank{#1}{\gappto\FrankVocBuf{#2\FrankChunkSep}}%
              {\gappto\FrankVocBuf{{#1#2}\FrankChunkSep}}}

% ------------------------- the word line -----------------------------
% lib/wordline.lua reads a chunk's words for the PDF.  It is loaded ONCE, here,
% with dofile: ordinary Lua in a file of its own, out of reach of every
% \directlua trap this file warns about, and held to the same cases as the
% Python and the JavaScript parsers by tests/fixtures/wordline.json.
\directlua{frank_wordline = dofile("\FrankLib/wordline.lua")}
% \FrankWords{fa}{words}: what wordline.lua's tex() writes -- one
% \FrankWord{text}{reading} per word and \FrankWordSep between two, the
% characters drawn from fa and only the division and the readings from the
% line.  A line that does not fit its text stops the build (tex.error names
% the chunk) rather than print a reading over the wrong characters.  Blank
% words are the text alone: the macro that carried them has said so.
\newcommand{\FrankWords}[2]{\ifblank{#2}{#1}{%
  \directlua{tex.sprint(frank_wordline.tex("\luaescapestring{#1}","\luaescapestring{#2}"))}}}
\newcommand{\FrankWord}[2]{\ifblank{#2}{#1}{\FrankRuby{#1}{#2}}}
% \FrankWordP{opening}{text}{reading}{closing}{hanging}: a word with the
% bracket that opens before it and the marks that close after it (「おい」,
% 行きました。), which no line may be broken away from.  #5 is the 、 。 ， ．
% at the very end: the only marks Japanese and Chinese setting lets hang past
% the edge of the line (burasagari), #4 whatever closes before them.
%   - All of it fits the measure: one box.
%   - It does not, but it would without #5: one box with #5 in no width,
%     hanging past the edge.  The long word and its comma at the end of the
%     chunk column (切ってみると、, a third of the line): the gutter to the
%     gloss, 0.04 of the line, holds the mark.  Two boxes with no break
%     between them ran 6.8pt into it with an overfull warning, twice, in the
%     fixture edition divided into words.
%   - Not even that: the reading macro is handed all of it, and its own
%     measured fallback takes over -- the reading on a line of its own and
%     the text wrapping beneath, what a chunk too wide for its column gets.
% A word with no reading is its marks and its text, broken the usual way.
\newcommand{\FrankWordP}[5]{%
  \ifblank{#3}{#1#2#4#5}{\FrankWordPRuby{#1}{#2}{#3}{#4}{#5}}}
% The reading may reach past its word over the marks beside it, which have
% no reading of their own: a pinyin wider than its characters (zǎoshang over
% 早上) otherwise pushed the comma a gap away from the word it closes.  Box 2
% is the word under its reading, \dimen2 how far the reading reaches past the
% word on either side, \dimen4/\dimen6 how much of that the opening and the
% closing marks take, never more than the marks are wide.
\newcommand{\FrankWordPRuby}[5]{%
  \begingroup
  \setbox2\hbox{\FrankRuby{#2}{#3}}%
  \setbox4\hbox{#2}%
  \dimen2=\dimexpr(\wd2-\wd4)/2\relax
  \ifdim\dimen2<0pt \dimen2=0pt \fi
  \setbox6\hbox{#1}%
  \ifdim\wd6<\dimen2 \dimen4=\wd6 \else\dimen4=\dimen2 \fi
  \setbox6\hbox{#4#5}%
  \ifdim\wd6<\dimen2 \dimen6=\wd6 \else\dimen6=\dimen2 \fi
  \setbox0\hbox{#1\kern-\dimen4 \copy2\kern-\dimen6 #4#5}%
  \ifdim\wd0>\hsize
    \setbox6\hbox{#4}%
    \ifdim\wd6<\dimen2 \dimen6=\wd6 \else\dimen6=\dimen2 \fi
    \setbox0\hbox{#1\kern-\dimen4 \copy2\kern-\dimen6 #4\rlap{#5}}%
    \ifdim\wd0>\hsize \FrankRuby{#1#2#4#5}{#3}\else\leavevmode\box0 \fi
  \else\leavevmode\box0 \fi
  \endgroup}
% Between two words: no width, a little stretch, and a legal break.  babel
% puts its glue between two characters and never beside a box, and a word
% with its ruby is an unbreakable box.
\newcommand{\FrankWordSep}{\hskip 0pt plus .1em\relax}
% A \chw or \chrw whose words are blank is the chunk without them, which has
% a name of its own.  Said once per chunk, as it is set.
\newcommand{\FrankNeedWords}[4]{% {\chw}{\ch}{fa}{words}
  \ifblank{#4}{\PackageError{frank-preamble}{%
    \string#1 has blank words for `#3':\MessageBreak
    use \string#2, which is this chunk without words}{%
    The last argument of \string#1 is the chunk's word line.}}{}}

% ------------------------- the reading pass --------------------------
% \FrankAloudAppend{kana}{tr}{fa}: a chunk's part in the reading pass, where
% the language has one.  The reading is the kana where the language has a
% reading and the transliteration where it has not (Chinese: the pinyin), and
% never taken from the words: a kanbun reading reorders them, and a Chinese tr
% writes the changed tones of 一 and 不.  \gappto: nothing is expanded until
% the pass is set.
\newcommand{\FrankAloudAppend}[3]{%
  \ifFrankHasAloud
    \ifFrankHasReading
      \gappto\FrankAloudBuf{\FrankAloud{#1}{#3}\FrankAloudSep}%
    \else
      \gappto\FrankAloudBuf{\FrankAloud{#2}{#3}\FrankAloudSep}%
    \fi
  \fi}
% \FrankAloud{reading}{fa}: wordline.lua's tex_aloud -- the reading in
% \FrankAloudText, then the closing punctuation of the text that the reading
% leaves out (やまへしばかりに over 山へ柴刈りに、 reads やまへしばかりに、); the text
% itself where the reading is blank, which is every \chp and a Japanese \ch.
\newcommand{\FrankAloud}[2]{%
  \directlua{tex.sprint(frank_wordline.tex_aloud("\luaescapestring{#1}","\luaescapestring{#2}"))}}
% The gap between two chunks of that pass: visible, because a reading run on
% without a break hides where one chunk ends, and the line's preferred break.
% Kana break between any two characters, and the fixture edition cut いきました
% after its い, stranded at a line's end after the gap.  -1 and no more: under
% \raggedright every line is as good as any other, so the least penalty tips
% the tie, and at -50 -- more than \linepenalty squared -- TeX broke at every
% chunk and set each on a line of its own.  A language file may set its own.
\providecommand{\FrankAloudSep}{\penalty-1\hskip .6em plus .2em\relax}

\newcommand{\ch}[5]{%
  \ifFrankCollect
    \FrankVocAppend{#1}{#2}%
    \xappto\FrankBuf{\Strip{#2}\FrankChunkSep}%
    \FrankAloudAppend{}{#3}{#2}%
  \else
  \ifFrankFlow
    #2\,\penalty300 \mbox{{\SzGloss\color{gloss}%
      \foreignlanguage{\FrankGlossName}{(\FrankRoman{\textit{#3}} — #5)}}}\ %
  \else
    \FrankChunkRow{\FrankBreaks{#2}}{\glbody{#3}{#4}{#5}}%
  \fi\fi}

% \chr{colour}{fa}{kana}{tr}{voc}{en} -- the chunk with its reading.  The
% buffers receive exactly what \ch would give them, except that the pass-1
% copy carries the ruby, so the reader's own attempt is made WITH the kana
% over each chunk (and the bare copy without).  \jruby is the language file's
% (ja.tex); elsewhere it prints the text alone and the kana line is not set,
% so for a language without a reading this is \ch with an argument ignored.
\newcommand{\chr}[6]{%
  \ifFrankCollect
    \FrankVocAppend{#1}{\jruby{#2}{#3}}%
    \xappto\FrankBuf{\Strip{#2}\FrankChunkSep}%
    \FrankAloudAppend{#3}{#4}{#2}%
  \else
  \ifFrankFlow
    \jruby{#2}{#3}\,\penalty300 \mbox{{\SzGloss\color{gloss}%
      \foreignlanguage{\FrankGlossName}{(\FrankRoman{\textit{#4}} — #6)}}}\ %
  \else
    \FrankChunkRow{\jruby{\FrankBreaks{#2}}{#3}}{\glbodyr{#3}{#4}{#5}{#6}}%
  \fi\fi}

% \chw{colour}{fa}{tr}{voc}{en}{words}, \chrw{colour}{fa}{kana}{tr}{voc}{en}{words}
% -- \ch and \chr line for line, except where the text is SET: pass 1 and the
% chunk column take \FrankWords, and the chunk's kana is no longer ruby over
% the whole chunk; it stays the first line of the gloss.  The pass-1 copy goes
% in through \FrankVocAppend's \gappto, unexpanded -- an \xappto of anything
% but the plain text is fatal, a ruby box being nothing an \edef can take
% apart -- and the bare copy is \Strip{fa} exactly as ever.
\newcommand{\chw}[6]{%
  \ifFrankCollect
    \FrankVocAppend{#1}{\FrankWords{#2}{#6}}%
    \xappto\FrankBuf{\Strip{#2}\FrankChunkSep}%
    \FrankAloudAppend{}{#3}{#2}%
  \else
  \FrankNeedWords{\chw}{\ch}{#2}{#6}%
  \ifFrankFlow
    \FrankWords{#2}{#6}\,\penalty300 \mbox{{\SzGloss\color{gloss}%
      \foreignlanguage{\FrankGlossName}{(\FrankRoman{\textit{#3}} — #5)}}}\ %
  \else
    \FrankChunkRow{\FrankWords{#2}{#6}}{\glbody{#3}{#4}{#5}}%
  \fi\fi}

\newcommand{\chrw}[7]{%
  \ifFrankCollect
    \FrankVocAppend{#1}{\FrankWords{#2}{#7}}%
    \xappto\FrankBuf{\Strip{#2}\FrankChunkSep}%
    \FrankAloudAppend{#3}{#4}{#2}%
  \else
  \FrankNeedWords{\chrw}{\chr}{#2}{#7}%
  \ifFrankFlow
    \FrankWords{#2}{#7}\,\penalty300 \mbox{{\SzGloss\color{gloss}%
      \foreignlanguage{\FrankGlossName}{(\FrankRoman{\textit{#4}} — #6)}}}\ %
  \else
    \FrankChunkRow{\FrankWords{#2}{#7}}{\glbodyr{#3}{#4}{#5}{#6}}%
  \fi\fi}

% a chunk that needs no gloss
% \chp{colour}{fa} -- the same, for a chunk that needs no gloss
\newcommand{\chp}[2]{%
  \ifFrankCollect
    \FrankVocAppend{#1}{#2}%
    \xappto\FrankBuf{\Strip{#2}\FrankChunkSep}%
    \FrankAloudAppend{}{}{#2}%
  \else
  \ifFrankFlow #2\ \else
    \par\noindent\hbox to\linewidth{\parbox[t]{\ChunkCol}{\raggedright \FrankBreaks{#2}}\hfil}%
    \par\addvspace{1.1ex}\fi\fi}

% The passes.  The body is read once silently to fill the buffers, so that
% pass 1 -- the whole sentence, with its help, unglossed -- can be set BEFORE
% the chunks.  That is the reader's own attempt, made before any help arrives.
% Then, where the chunks carry words, the reading alone; then the chunks; then,
% where the language has them, the bare text and the last pass -- the
% alternate face, or the vertical setting for a language that has one (the
% vertical pass takes the place of the font pass, so no language has both,
% and five is the most passes there are).
\NewEnviron{frank}{%
  \par\FrankBufReset
  {\FrankCollecttrue\BODY}%
  \FrankVocalPass
  \ifFrankHasAloud\FrankAloudPass\fi
  \begingroup\SzChunk\color{ink}
  \begin{otherlanguage}{\FrankLangName}\raggedright
  \BODY
  \par\end{otherlanguage}\endgroup
  \ifFrankHasBare\FrankPlainPass\fi
  \ifFrankHasVert\FrankVertPass\else\ifFrankHasAlt\FrankNastPass\fi\fi
  \par\addvspace{3.4ex}%
}

% The paragraph's number, set small and grey at the language's margin, just
% above the reader's own first attempt at it.  Printed as given: assemble.py
% writes it in the language's digits.
\newcommand{\parnum}[1]{%
  \par\addvspace{1.4ex}%
  {\SzGloss\color{faint}
   \begin{otherlanguage}{\FrankLangName}\raggedright\mbox{#1}\par\end{otherlanguage}}%
  \par\addvspace{0.2ex}}

\newcommand{\FrankVocalPass}{%
  \par\addvspace{0.6ex}%
  {\SzPlain\color{ink}
   \begin{otherlanguage}{\FrankLangName}\raggedright\FrankVocBuf\par\end{otherlanguage}}%
  \par\addvspace{2.4ex}%
  {\centering\color{faint}\rule{0.18\linewidth}{0.4pt}\par}%
  \par\addvspace{1.8ex}}

% the READING pass: every chunk's reading alone, a gap between two chunks, set
% as the plain pass is and closed with pass 1's short rule -- the sound of the
% sentence, still before any meaning.  The last chunk's gap comes off first.
\newcommand{\FrankAloudPass}{%
  {\SzPlain\color{ink}
   \begin{otherlanguage}{\FrankLangName}\raggedright
     \FrankAloudBuf\unskip\unpenalty\par\end{otherlanguage}}%
  \par\addvspace{2.4ex}%
  {\centering\color{faint}\rule{0.18\linewidth}{0.4pt}\par}%
  \par\addvspace{1.8ex}}

\newcommand{\FrankPlainPass}{%
  \par\addvspace{2.2ex}%
  {\SzPlain\color{ink}
   \begin{otherlanguage}{\FrankLangName}\raggedright\FrankBuf\par\end{otherlanguage}}}

% the FONT pass: the bare text in the alternate face (nastaliq, for Persian)
\newcommand{\FrankNastPass}{%
  \par\addvspace{1.6ex}%
  {\FrankAltFont\SzNast\color{ink}
   \begin{otherlanguage}{\FrankLangName}%
     \raggedright\tolerance=3000\emergencystretch=3em
     \FrankBuf\par
   \end{otherlanguage}}}

% ------------------- structural marks -------------------------------
% A page turn happens only where the original paragraph ends, so that the
% bare and fourth passes are never visible beside their own glosses.
% \parend closes a paragraph, \chapend closes a chapter -- the sun is the
% mark the 1951 Sepehr edition itself uses between sections.
\newcommand{\parend}{%
  \par\addvspace{2.4ex}%
  {\centering\color{faint}\normalfont\large\P\par}%
  \newpage}

\newcommand{\chapend}{%
  \par\addvspace{2.6ex}%
  {\centering\color{faint}\symbolfont\large ☼\,☼\,☼\par}%
  \newpage}

\newcommand{\pline}[2][\SzPlain]{%
  {#1\color{ink}\begin{otherlanguage}{\FrankLangName}\centering #2\par\end{otherlanguage}}}

\newcommand{\numbersep}[1]{%
  \par\addvspace{2ex}{\centering\color{faint}\rule{0.32\linewidth}{0.4pt}\par}%
  \pline[\SzPlain]{#1}\addvspace{1.6ex}}

% THE CHAPTER'S NAME, under the number \chapopen has just printed.  Optional:
% a chapter without one is written exactly as every chapter was before names
% existed, and both the PDF and the reader fall back to the number alone.
% It names; it does not number.  The number is \chapopen's and is printed
% already, so nothing here is counted and nothing downstream can be shifted
% by adding, changing or removing a name.
\newcommand{\chapname}[1]{%
  \par\vspace{-4mm}%
  {\fontsize{15}{24}\selectfont\color{ink}
   \begin{otherlanguage}{\FrankLangName}\centering #1\par\end{otherlanguage}}%
  \vspace{6mm}\par}

% A SECTION INSIDE A CHAPTER: a place between two paragraphs, and nothing
% more.  It has NO NUMBER OF ITS OWN and no effect whatever on the numbering
% of the paragraphs, which is \parnum's alone -- the same guarantee \parstart
% carries, for the same reason: verify_book.py proves the body still
% reproduces the source word for word, and a mark between two paragraphs
% must not move a single one of them.
%
% Smaller than \chapopen's number, as a subheading is smaller than a heading,
% and on the shorter centred rule so the two cannot be mistaken for each
% other at a glance.  A whole-line comment above it in the file is consumed
% with its newline and cannot reach the page (lib/timestamp.py).
\newcommand{\secmark}[1]{%
  \par\addvspace{3.4ex}%
  {\centering\color{faint}\rule{0.18\linewidth}{0.4pt}\par}%
  \vspace{2.5mm}%
  {\fontsize{15}{24}\selectfont\color{ink}
   \begin{otherlanguage}{\FrankLangName}\centering #1\par\end{otherlanguage}}%
  \addvspace{2.8ex}\par}

% opens a chapter: the number, in the alternate face where the language has
% one (nastaliq, gothic) and in the main face otherwise, on a fresh page
\newcommand{\chapopen}[1]{%
  \chapstart{#1}%
  \thispagestyle{fancy}%
  \vspace*{6mm}%
  {\centering\color{faint}\rule{0.30\linewidth}{0.4pt}\par}%
  \vspace{4mm}%
  {\ifFrankHasAlt\FrankAltFont\fi\fontsize{22}{34}\selectfont\color{ink}
   \begin{otherlanguage}{\FrankLangName}\centering #1\par\end{otherlanguage}}%
  \vspace{5mm}%
  {\centering\color{faint}\rule{0.30\linewidth}{0.4pt}\par}%
  \vspace{7mm}\par}
